% Generated by Claude Opus 5 (Anthropic) --- 2026-09-10
\section{Roadmap}
\label{sec:roadmap}

The anomaly of Table~\ref{tab:occgrid}---a callback whose typical cost falls
under tracing while its worst case grows fivefold---needs its root cause
established and mitigated. On the model side we aim to extend FISH towards
ARM-based SoC devices and extract task models out of functioning cyber-physical
systems; LTTng, ROS~2 and \texttt{nsys} are all available there, so this is more
labour than research. GPU profiling should become lighter: injecting CUPTI in
place of \texttt{nsys} would collect only the events the model needs, and the
same injection mechanism generalises to AMD through ROCm and to ARM Mali.

The host layer exists for deployments with more than one actor: several hosts, or
several containerised applications that communicate. The model generalises to
them unchanged. The host is the only layer whose cardinality is one today, and a
second one is a sibling subtree, since the hierarchy asks only that the children
of a layer partition the layer below. The rule that lays a $\mathit{msg}$ edge
wherever a publisher and a subscription share a topic never asks whether the two
sit under the same host, so traffic between hosts, or between separately
containerised applications, is recovered by the same construction. What the
setting adds is a common time base: an interaction measured across hosts is worth
no more than the offset estimated between their clocks.

FISH produces task graphs, which is half of the envisioned pipeline. From here we
move to the formalisation of task-set generation, over selected subsets of the
model or over its entirety, and to the analysis that consumes it. The
schedule-abstraction graph~\cite{nasri-2017-sagproposed} is a candidate: a
callback that submits GPU work and blocks until it completes is a self-suspending
task, and the analysis has been extended to that
class~\cite{nidhi-2024-sagedd}.
